\subsection{Benchmark scope and experimental units}
The development process separated geometries rather than individual routes or
random seeds. The principal training screen uses nine difficult instance
variants: a280, gil262, gr229, kroA200, pr226, rd400, si175, ts225 and tsp225.
Each geometry contributes one fixed variant, whose mandatory, optional and
demand sets and prize quota are retained unchanged across methods. Complete
instance identifiers and input hashes are stored in the batch manifests.
For these nine training files, the local SHA-256 digests match the corresponding
public file-information records in the Mendeley version-2 deposit cited by
the earlier PCCTP work~\cite{climaco2021data}. This identity check covers the
nine selected input variants; the remaining benchmark files and historical
evaluation subsets are outside its verified scope.
These selected cases are not a representative random sample of all PCCTP
instances and do not reproduce the full published benchmark evaluation.

The completed 60-second validation screen uses seven additional geometries.
Those geometries had already informed an earlier configuration and are
therefore treated as reused validation data, rather than untouched
confirmation. Twelve further test geometries remain unopened. Changing the
time budget or random seeds does not restore independence to data previously
used during method development. Results from selected training instances,
reused validation and conditional component replays are reported separately.

An inspection of the nine training inputs also distinguishes coverage
requirements from prize requirements. In gr229, kroA200, pr226 and rd400, the
mandatory locations already cover every demand. Coverage constraints are
consequently redundant in these four complete instances, although the prize
quota and routing decisions remain restrictive. An improvement on these
cases cannot be attributed to restoring uncovered demands. This structural
observation follows directly from the inputs and does not identify the cause
of any observed cost difference.

\subsection{Baseline and computing environment}
The reference implementation is the supplied MDM-GRASP associated with
the data-mining study~\cite{climaco2026mdm}. It is compiled separately from
the experimental Python drivers. Two isolated corrections initialize an
otherwise uninitialized decoder index and bound elite-bank access by the
number of occupied entries. The legacy source is preserved. Results of the
corrected implementation are identified accordingly; they are not described
as measurements of an unmodified published executable.

The reference configuration uses 70 maximum iterations, construction
parameter $\alpha=0.6$, an elite-bank capacity of ten, absolute mining support
two, and a stability threshold of five percent of the iteration limit.
Absolute support two corresponds to 20 percent only when the bank contains
ten solutions; the experiment does not silently recompute it for a partially
occupied bank. The reference policy bypasses experimental learned rankings
and adaptive timing decisions. Instrumentation and collection still impose
overhead and are not assumed to be cost-free.

The local environment is an Apple M2 with 16 GiB of memory and eight physical
cores. The recorded native compiler is GCC 15.1.0; the research driver uses
Python 3.9, SciPy 1.13 and HiGHS through the available Python interfaces.
The ordering component is LKH 3.0.6. Program and binary hashes identify the
executables used by each batch. Only one research optimizer is scheduled at
a time. This scheduling constraint does not establish a constant background
load, temperature or effective processor frequency. Raw seconds are not
converted into speedups against the different processor used in a prior
paper by scaling clock frequencies.

\subsection{Completed experiments and longer-horizon design}
The completed total-time training experiment allocates 60 seconds to each
method, using nine geometries and three matched seeds. The completed
validation screen uses seven geometries and five seeds. Both start from a
fresh, charged fallback and include initialization within the outer budget.
The fixed-support replay instead starts from the saved pre-repair cycle and
uses its recorded remaining allowance. Generation of that shared state is
excluded from the replay, which is therefore a conditional component test.
The matched-destruction MIP experiment similarly shares saved MDM states;
its comparisons must not be interpreted as complete algorithms with free
incumbent generation.

The longer-horizon training design allocates 300 seconds to each of three
methods: the reference, a hybrid with ten seconds of reference initialization,
and a hybrid with 60 seconds of reference initialization. All use the same
nine variants and seeds 70--72. Method order is shuffled reproducibly within
each instance--seed block. Destruction, repair, ordering and acceptance
parameters are frozen before collection. Both hybrid configurations are
retained in the report, regardless of their relative performance. This
design schedules 81 executions and remains an exploratory training screen.
Its incomplete outcomes are excluded from the completed-results table.

A separate prospective design addresses the 70-iteration stopping criterion.
An initial nine-case calibration with seed 80 must complete and pass its
audits before the ten-seed comparison starts. In that subsequent comparison,
each of the nine training variants uses seeds 81--90. A reference run first
completes all 70 iterations and provides its measured process-wall duration
$T_{g,s}$. Each hybrid then receives that same outer solution-eligibility
horizon, including its own fallback, setup and initialization. Initialization
is capped at $\min\{10,T_{g,s}\}$ or $\min\{60,T_{g,s}\}$, respectively.
Neither the reference route nor its objective value is supplied as a hybrid
starting solution. Independent post-exit auditing does not enlarge a search
budget.

The resulting horizon is conditioned on the measured reference duration;
it is not an independently fixed common deadline. The reference necessarily
precedes the candidates in a pair. Shuffling pair order and the two candidate
orders does not remove that chronological asymmetry, which is disclosed as
a limitation. The fixed-300-second design supplies a complementary comparison
with shuffled order among all three methods. The ten-seed design schedules
90 reference runs and 180 candidate runs if every reference completes. It
is not yet a completed experiment and does not constitute a replication of
the full published instance set.

\subsection{Outcome aggregation and failure handling}
For geometry $g$, let $\bar c_{B,g}$ and $\bar c_{H,g}$ be the arithmetic
mean final costs over all its scheduled seeds for the reference and a
candidate, respectively. The primary descriptive effect is
\begin{equation}
 G_g=100\frac{\bar c_{B,g}-\bar c_{H,g}}{\bar c_{B,g}},
 \qquad
 \bar G=\frac{1}{|\mathcal G|}\sum_{g\in\mathcal G}G_g.
 \label{eq:geometry-effect}
\end{equation}
Positive values favor the candidate. This equally weighted geometry average
differs from averaging seed-level percentages or pooling objective values
across differently scaled instances. Seeds describe variability conditional
on an instance; they do not multiply the number of independent geometries.
All scheduled configurations, ties, losses and failed runs remain visible.
Best-seed selection is not used for the primary comparison.

The nine-case training screen requires a mean reduction of at least
$0.2\%$, at least six positive geometry means, at most one negative mean,
and no missing or invalid scheduled comparisons. These thresholds are an
exploratory investment rule, not a significance test or an acceptance
criterion for publication. The earlier validation screen retains its
original, different threshold and its observed failure. It is not relabeled
as successful because a later training experiment looks favorable.

A safety timeout of 43,200 seconds in an iteration-limited reference run
means censoring, not completion. Without a completed reference, no valid
$T_{g,s}$ is available for either corresponding candidate. Both candidate
slots remain unavailable in the planned denominator. A last checkpoint is
retained as an observation but does not replace the missing 70-iteration
endpoint. No complete cost-only primary result is formed by dropping such
pairs. Infrastructure interruptions are recorded separately from genuine
algorithm or solver failures.

\subsection{Deadlines, suspension and independent audits}
Solution eligibility is determined from absolute monotonic event timestamps
relative to the parent start time. A solution must be observed by the outer
deadline and pass validation against the original instance. Cleanup or process return after that deadline does
not create additional search time, nor does it erase an already eligible
incumbent. Actual process durations, safety timeouts and cleanup overruns
are retained separately. Descriptive curves at 1, 10, 30, 60, 120 and
300 seconds do not replace the primary 300-second endpoint.

The first longer-horizon collection was interrupted after host suspension
was identified in the operating-system power log. Its 44 completed records
are preserved with an explicit environmental-interruption manifest. They
are not treated as a controlled time comparison. The replacement collection
repeats all 81 scheduled executions in a separate directory, rather than
combining apparently unaffected early cases with favorable later repetitions.
It retains the same instances, seeds, parameters, method ordering and binary
versions.

During the replacement collection, a monitor samples monotonic and wall
clocks once per second. A monitoring gap above five seconds, a difference
between clock increments above three seconds, or a recorded system-sleep
event invalidates the block and stops the queue. Temporary idle-sleep
prevention does not prevent sleep caused by closing the laptop lid.
Each block records environmental diagnostics and hashes of its three
results. A separate auditor checks those records. These operational checks
screen for suspension and substantial scheduling interruptions; they do not
prove identical CPU or thermal conditions. Equivalent protection is required
before accepting the prospective iteration-limited comparisons.

An independent scalar reader checks mandatory visits, unique locations,
coverage, prize collection and the cycle cost against the original input
matrix. Audits also reconstruct eligible incumbent selection and check
input, source, binary and result hashes. Feasibility verification and a
successful environment check address different questions; neither substitutes
for the other. Historical finite verification records without source hashes
are retained as historical evidence and cannot by themselves certify the
final source revision. This partial manuscript makes no confirmatory
significance claim or new best-known-solution claim from the exploratory
comparisons.
